% ---------------------------------------------------------------------------
% ACT A -- MOTOR-IN-DUCT THERMAL MAP -- one-page customer-facing result sheet.
% Written under the DEMO STANDARD v2 output rules R1-R10.
% Every user-visible number on this sheet traces to a source artifact by file
% and line. That mapping lives in README_SOURCES.md in this directory; it is
% not repeated here, so that this file carries no internal identifiers at all.
% Self-contained: no external figure file, no packages beyond a base TeX Live
% install (no tikz/pgfplots on this box). The plot is a native LaTeX picture,
% so the output is vector at any zoom.
% Compile: pdflatex ACT_A_thermal_map_sheet.tex
% ---------------------------------------------------------------------------
\documentclass[9pt,a4paper]{article}
\usepackage[margin=10mm,top=7mm,bottom=5mm]{geometry}
\usepackage{booktabs}
\usepackage{array}
\usepackage{xcolor}
\usepackage{graphicx}
\usepackage[T1]{fontenc}

\definecolor{safe}{rgb}{0.88,0.94,0.88}
\definecolor{limit}{rgb}{0.72,0.12,0.12}
\definecolor{rulegrey}{rgb}{0.45,0.45,0.45}
\definecolor{boxgrey}{rgb}{0.96,0.96,0.94}

\renewcommand{\footnotesize}{\fontsize{6.9}{7.9}\selectfont}
\renewcommand{\scriptsize}{\fontsize{6.4}{7.5}\selectfont}
\renewcommand{\arraystretch}{1.02}
\setlength{\parindent}{0pt}
\setlength{\parskip}{0.5pt}
\pagestyle{empty}

\newcommand{\rolesig}[1]{\textbf{\footnotesize#1}}

\begin{document}
\fontsize{8}{9.5}\selectfont
\enlargethispage{30pt}

% ---------------------------------------------------------------- title -----
{\fontsize{14}{16}\selectfont\bfseries Motor in a cooling duct: thermal envelope}\\[1.5pt]
{\fontsize{8.8}{10}\selectfont Peak housing temperature over 16 solved operating
points --- 4 heat loads $\times$ 4 airspeeds}\\[2pt]
{\color{rulegrey}\rule{\textwidth}{0.9pt}}\\[1pt]
{\footnotesize\textbf{Solver:} OpenFOAM \texttt{chtMultiRegionSimpleFoam},
steady conjugate heat transfer, $k$--$\omega$ SST.
Axisymmetric duct sector. Cooling air at $288\,\mathrm{K}$
($14.85\,^{\circ}\mathrm{C}$). Housing temperature limit
$200\,^{\circ}\mathrm{C}$ (customer-stated).}

% ---------------------------------------------------------------- table 1 ---
\vspace{2pt}
\textbf{Table 1 --- Peak housing temperature $T_{\max}$, $^{\circ}\mathrm{C}$.}
Sixteen solved points; all values read from the solution at full write
precision.

\vspace{2pt}
\begin{tabular}{r rrrr l}
\toprule
 & \multicolumn{4}{c}{\textbf{Airspeed $U_{\infty}$, m\,s$^{-1}$}} & \\
\cmidrule(lr){2-5}
\textbf{Heat load $P$, W} & \textbf{10} & \textbf{20} & \textbf{30} & \textbf{40}
 & \textbf{Numerical uncertainty on every value in the row} \\
\midrule
\textbf{80}  &  38.1374 &  29.0795 & 25.5462 & 23.5897 & not quantified (single mesh level) \\
\textbf{155} &  59.9609 &  42.3896 & 35.5104 & 31.6747 & not quantified (single mesh level) \\
\textbf{230} &  81.7844 &  55.6997 & 45.4746 & 39.7598 & not quantified (single mesh level) \\
\textbf{305} & 103.6078 &  69.0098 & 55.4389 & 47.8448 & not quantified (single mesh level) \\
\bottomrule
\end{tabular}

\vspace{1.5pt}
{\footnotesize Every point sits on one mesh level. An error bar needs the same
point on two further meshes. None is quoted for these sixteen numbers.}

\vspace{2pt}
\textbf{Table 2 --- Margin to the $200\,^{\circ}\mathrm{C}$ limit, at the
worst airspeed $10\,\mathrm{m\,s^{-1}}$.} Margin is the limit minus the
solved peak.

\vspace{2pt}
{\fontsize{7.2}{8.5}\selectfont
\begin{tabular}{r r r l}
\toprule
Heat load $P$, W & Peak $T_{\max}$, $^{\circ}\mathrm{C}$ &
Margin, K & Reading \\
\midrule
80  &  38.1374 & 161.8626 & clears the limit \\
155 &  59.9609 & 140.0391 & clears the limit \\
230 &  81.7844 & 118.2156 & clears the limit \\
\textbf{305} & \textbf{103.6078} & \textbf{96.3922} & worst solved point \\
\bottomrule
\end{tabular}}

\vspace{1.5pt}
{\footnotesize The coolest point is $23.5897\,^{\circ}\mathrm{C}$, at
$80\,\mathrm{W}$ and $40\,\mathrm{m\,s^{-1}}$. Margins carry no error bar,
for the same reason as Table 1.}

% ------------------------------------------------- plot + assumption beat ---
\vspace{1.5pt}
\begin{minipage}[t]{0.485\textwidth}
\textbf{Figure 1 --- Operating envelope.} Peak housing temperature against
airspeed, one curve per heat load. Shaded band is below the limit.

\vspace{2pt}
\setlength{\unitlength}{1mm}
\begin{picture}(88,63)(-8,-8)
  % ---- safe region: everything below the 200 degC line -------------------
  \put(0,0){{\color{safe}\rule{84mm}{53.33mm}}}
  % ---- axes ---------------------------------------------------------------
  \linethickness{0.4pt}
  \put(0,0){\line(1,0){84}}
  \put(0,0){\line(0,1){57}}
  % ---- y ticks and labels, degC ------------------------------------------
  \multiput(0,13.33)(0,13.33){4}{\line(-1,0){1.4}}
  \put(-2.2,-0.9){\makebox(0,0)[r]{\scriptsize 0}}
  \put(-2.2,12.43){\makebox(0,0)[r]{\scriptsize 50}}
  \put(-2.2,25.77){\makebox(0,0)[r]{\scriptsize 100}}
  \put(-2.2,39.10){\makebox(0,0)[r]{\scriptsize 150}}
  \put(-2.2,52.43){\makebox(0,0)[r]{\scriptsize 200}}
  \put(-7.0,30){\makebox(0,0)[c]{\rotatebox{90}{\scriptsize $T_{\max}$, $^{\circ}\mathrm{C}$}}}
  % ---- x ticks and labels, m/s -------------------------------------------
  \multiput(28,0)(28,0){3}{\line(0,-1){1.4}}
  \put(0,-4.2){\makebox(0,0)[c]{\scriptsize 10}}
  \put(28,-4.2){\makebox(0,0)[c]{\scriptsize 20}}
  \put(56,-4.2){\makebox(0,0)[c]{\scriptsize 30}}
  \put(84,-4.2){\makebox(0,0)[c]{\scriptsize 40}}
  \put(42,-7.6){\makebox(0,0)[c]{\scriptsize Airspeed $U_{\infty}$, m\,s$^{-1}$}}
  % ---- the 200 degC limit line -------------------------------------------
  \linethickness{0.7pt}
  \multiput(0,53.33)(4,0){21}{{\color{limit}\rule{2.2mm}{0.6pt}}}
  \put(2,55.0){\makebox(0,0)[l]{{\color{limit}\scriptsize
      $200\,^{\circ}\mathrm{C}$ housing limit (customer-stated)}}}
  \put(2,49.3){\makebox(0,0)[l]{\scriptsize safe region}}
  % ---- curves: 4 measured samples joined by straight segments -------------
  \linethickness{0.5pt}
  % 80 W  -> 10.17 7.75 6.81 6.29
  \qbezier(0,10.17)(14,8.96)(28,7.75)
  \qbezier(28,7.75)(42,7.28)(56,6.81)
  \qbezier(56,6.81)(70,6.55)(84,6.29)
  \put(0,10.17){\circle*{1.1}}\put(28,7.75){\circle*{1.1}}
  \put(56,6.81){\circle*{1.1}}\put(84,6.29){\circle*{1.1}}
  % 155 W -> 15.99 11.30 9.47 8.45
  \qbezier(0,15.99)(14,13.65)(28,11.30)
  \qbezier(28,11.30)(42,10.39)(56,9.47)
  \qbezier(56,9.47)(70,8.96)(84,8.45)
  \put(0,15.99){\circle*{1.1}}\put(28,11.30){\circle*{1.1}}
  \put(56,9.47){\circle*{1.1}}\put(84,8.45){\circle*{1.1}}
  % 230 W -> 21.81 14.85 12.13 10.60
  \qbezier(0,21.81)(14,18.33)(28,14.85)
  \qbezier(28,14.85)(42,13.49)(56,12.13)
  \qbezier(56,12.13)(70,11.37)(84,10.60)
  \put(0,21.81){\circle*{1.1}}\put(28,14.85){\circle*{1.1}}
  \put(56,12.13){\circle*{1.1}}\put(84,10.60){\circle*{1.1}}
  % 305 W -> 27.63 18.40 14.78 12.76
  \linethickness{0.9pt}
  \qbezier(0,27.63)(14,23.02)(28,18.40)
  \qbezier(28,18.40)(42,16.59)(56,14.78)
  \qbezier(56,14.78)(70,13.77)(84,12.76)
  \put(0,27.63){\circle*{1.3}}\put(28,18.40){\circle*{1.3}}
  \put(56,14.78){\circle*{1.3}}\put(84,12.76){\circle*{1.3}}
  % ---- legend, inside the axes where there is white space ----------------
  \put(55,45.5){\makebox(0,0)[l]{\scriptsize Heat load $P$:}}
  \linethickness{0.9pt}\qbezier(55,41.5)(58,41.5)(61,41.5)
  \put(58,41.5){\circle*{1.3}}
  \put(62.5,41.5){\makebox(0,0)[l]{\scriptsize 305\,W}}
  \linethickness{0.5pt}
  \qbezier(55,37.5)(58,37.5)(61,37.5)\put(58,37.5){\circle*{1.1}}
  \put(62.5,37.5){\makebox(0,0)[l]{\scriptsize 230\,W}}
  \qbezier(55,33.5)(58,33.5)(61,33.5)\put(58,33.5){\circle*{1.1}}
  \put(62.5,33.5){\makebox(0,0)[l]{\scriptsize 155\,W}}
  \qbezier(55,29.5)(58,29.5)(61,29.5)\put(58,29.5){\circle*{1.1}}
  \put(62.5,29.5){\makebox(0,0)[l]{\scriptsize 80\,W}}
\end{picture}

\vspace{1pt}
{\footnotesize Markers are the solved points of Table 1. Straight segments join
them; nothing between the four airspeeds is claimed. No uncertainty band is
drawn, because none exists at a single mesh level.}
\end{minipage}\hfill
\begin{minipage}[t]{0.485\textwidth}
\textbf{Table 3 --- Your resistance estimate against the solved answer, at
$305\,\mathrm{W}$.} Temperatures in $^{\circ}\mathrm{C}$. The ratio is taken on
the temperature \emph{rise} above the $288\,\mathrm{K}$ air, which is common to
both.

\vspace{2pt}
{\fontsize{7.2}{8.5}\selectfont
\begin{tabular}{r rr r rr}
\toprule
$U_{\infty}$ & \multicolumn{2}{c}{Hand estimate, $^{\circ}\mathrm{C}$} &
Solved & \multicolumn{2}{c}{Overprediction, $\times$} \\
\cmidrule(lr){2-3}\cmidrule(lr){5-6}
m\,s$^{-1}$ & duct corr. & flat plate & $^{\circ}\mathrm{C}$ & duct & flat plate \\
\midrule
10 & 329.1 & 195.1 & 103.6078 & \textbf{3.541} & \textbf{2.031} \\
20 & 197.3 & 120.3 &  69.0098 & 3.369 & 1.947 \\
30 & 148.0 &  92.4 &  55.4389 & 3.280 & 1.911 \\
40 & 121.6 &  77.4 &  47.8448 & \textbf{3.235} & \textbf{1.896} \\
\bottomrule
\end{tabular}}

\vspace{2pt}
{\footnotesize\textbf{The lumped resistance estimate is high by
$1.90\times$ to $3.54\times$.} Both variants overshoot. The overshoot exceeds
the $1.763\times$ spread between them, so the correlation choice is not the
cause. \emph{Why:} the hand method takes one resistance from a
fully-developed correlation. The real path is a short annular gap with
developing flow, and it removes heat better than that correlation allows.}

\vspace{2pt}
\textbf{Table 4 --- What the solved result says you could actually run at
$20\,\mathrm{m\,s^{-1}}$.}

\vspace{2pt}
{\fontsize{7.2}{8.5}\selectfont
\begin{tabular}{l r l}
\toprule
Reading & Heat load, W & Basis \\
\midrule
Solved and demonstrated        & \textbf{305} & solved \\
Implied to reach $120\,^{\circ}\mathrm{C}$ & 592 & extrapolated \\
Implied to reach $200\,^{\circ}\mathrm{C}$ & 1043 & extrapolated \\
\bottomrule
\end{tabular}}

\vspace{1.5pt}
{\footnotesize The two implied loads are arithmetic beyond the highest load ever
solved ($305\,\mathrm{W}$), on a proportionality checked only inside
$80$--$305\,\mathrm{W}$. \textbf{They are not measurements.} The demonstrated
statement is the first row: at least $305\,\mathrm{W}$ at every airspeed tested,
with at least $96.4\,\mathrm{K}$ of margin at the worst point.}
\end{minipage}

% -------------------------------------------------------- verification ------
\vspace{1.5pt}
{\color{rulegrey}\rule{\textwidth}{0.5pt}}\\[1pt]
\textbf{What is checked, and what is not}

\vspace{1pt}
\begin{minipage}[t]{0.60\textwidth}
\vspace{0pt}
{\fontsize{7.2}{8.5}\selectfont
\begin{tabular}{l l}
\toprule
\textbf{Check} & \textbf{Result} \\
\midrule
Acceptance criteria, fixed in advance & 16 of 16 points meet them \\
Trend with heat load, at fixed airspeed & rises in proportion, no reversal \\
Trend with airspeed, at fixed heat load & falls steeply, no reversal \\
Grid-refinement check & not available, none claimed \\
Energy-conservation check & not available, none claimed \\
\bottomrule
\end{tabular}}

\vspace{1.5pt}
{\footnotesize Two checks are absent, and neither is claimed. Both rows above
say so in the table rather than in prose.}
\end{minipage}\hfill
\begin{minipage}[t]{0.375\textwidth}
\vspace{0pt}
\colorbox{boxgrey}{\begin{minipage}{0.955\textwidth}
{\footnotesize\textbf{Read this before you use the numbers}\\[1pt]
$\bullet$ One mesh level per point, so there is no numerical error bar on any
temperature in Table 1.\\
$\bullet$ Thermal radiation is not modelled. No bound on its effect on these
temperatures is available.\\
$\bullet$ Buoyancy is switched off. Whether forced airflow genuinely dominates
is not established here, and the low-airspeed column is where that assumption
is least comfortable.\\
$\bullet$ Material properties are representative values, not measured for your
parts.\\
$\bullet$ The duct is treated as an axisymmetric sector; no circumferential
features are represented.\\
$\bullet$ Near-wall mesh resolution meets its target at $10$ and
$20\,\mathrm{m\,s^{-1}}$, and exceeds it at $30$ and $40\,\mathrm{m\,s^{-1}}$.
That limits what may be said about surface heat-transfer coefficients. The
temperature bound above is unaffected.\\
$\bullet$ Nothing is claimed between the four airspeeds solved.}
\end{minipage}}
\end{minipage}

% -------------------------------------------------------------- cost --------
\vspace{2pt}
{\footnotesize\textbf{Compute.} Sixteen operating points on this geometry.
Cost $\$0.49$, derived at $\$0.0513$ per core-hour. The rate is stated by the
machine's owner rather than metered here.}

% ------------------------------------------------------------- voices -------
\vspace{2pt}
{\color{rulegrey}\rule{\textwidth}{0.5pt}}

\vspace{1pt}
{\footnotesize
\rolesig{Lead Researcher.} The result worth your attention is the correction to
the sizing method, not the map. A lumped resistance estimate of this duct runs
high by two to three and a half times at every airspeed, worst at the low-speed
end.

\rolesig{Lead Engineer.} At $305\,\mathrm{W}$ and $10\,\mathrm{m\,s^{-1}}$ the
housing sits $96.4\,\mathrm{K}$ below your $200\,^{\circ}\mathrm{C}$ limit, and
the whole $80$--$305\,\mathrm{W}$ range clears it at all four airspeeds. Scale
by proportion within that range; do not read between the airspeeds, and do not
treat Table 4's extrapolated loads as a rating.

\rolesig{Lead Numericist.} The uncertainty column of Table 1 is empty on
purpose: one mesh level supports no error bar, and none is invented. A
quantified numerical uncertainty here is a mesh-refinement study.}

\vspace{1pt}
{\color{rulegrey}\rule{\textwidth}{0.4pt}}\\[0.5pt]
{\scriptsize 16 operating points solved on this geometry, 39,680 cells.}

\end{document}
